\documentclass{article}
\usepackage{amsmath,amssymb,amsthm}
\usepackage{algorithm,algorithmic}
\usepackage{fullpage}
\newtheorem{theorem}{Theorem}
\newtheorem{lemma}{Lemma}
\newtheorem{assumption}{Assumption}
\begin{document}

%=====================================================================
\section*{Algorithm Name}
Stochastic Gradient Descent with Warm Restarts (SGDR)
%=====================================================================

%=====================================================================
\section*{Mathematical Formulation}
Let $f:\mathbb{R}^d\rightarrow\mathbb{R}$ be the (possibly non‑convex) loss function.
At iteration $t\in\{0,1,\dots\}$ we draw a random sample $z_t$ and compute a stochastic
gradient 
\[
g_t:=\nabla f(w_t;z_t)\in\mathbb{R}^d .
\]

The learning rate $\eta_t$ follows a \emph{cosine annealing with warm restart}
schedule.  Fix a base period $M\in\mathbb{N}$ and a maximal step size
$\eta_{\max}>0$.  Define the \emph{cycle index}
\[
c(t):=\left\lfloor\frac{t}{M}\right\rfloor,\qquad
\theta(t):=\frac{t-Mc(t)}{M}\in[0,1).
\]
Then the learning rate at time $t$ is
\[
\boxed{\;
\eta_t \;=\; \eta_{\max}\,
\frac{1}{2}\Bigl(1+\cos\bigl(\pi\theta(t)\bigr)\Bigr)
\;}
\tag{1}
\]
i.e. it decays cosine‑wise from $\eta_{\max}$ to $0$ within each cycle of length $M$,
and is reset to $\eta_{\max}$ at the beginning of a new cycle.

The parameter update is the standard SGD step
\[
\boxed{\;
w_{t+1}=w_t-\eta_t\,g_t
\;}
\tag{2}
\]

Hyper‑parameters:
\begin{itemize}
  \item $\eta_{\max}$ – maximal step size,
  \item $M$ – length of a restart cycle,
  \item (optional) $T$ – total number of iterations.
\end{itemize}
%=====================================================================

%=====================================================================
\section*{Pseudocode}
\begin{algorithm}[H]
\caption{SGDR}
\begin{algorithmic}[1]
\REQUIRE Initial point $w_0\in\mathbb{R}^d$, maximal step size $\eta_{\max}$,
        cycle length $M$, total iterations $T$.
\FOR{$t=0$ \TO $T-1$}
    \STATE Compute cycle index $c\leftarrow\lfloor t/M\rfloor$;
    \STATE Compute normalized position $\theta\leftarrow (t-Mc)/M$;
    \STATE Set learning rate $\displaystyle\eta_t\leftarrow
          \frac{\eta_{\max}}{2}\bigl(1+\cos(\pi\theta)\bigr)$; \COMMENT{(1)}
    \STATE Sample $z_t$ and evaluate stochastic gradient $g_t=\nabla f(w_t;z_t)$;
    \STATE Update parameters $\displaystyle w_{t+1}\leftarrow w_t-\eta_t g_t$; \COMMENT{(2)}
\ENDFOR
\RETURN $w_T$
\end{algorithmic}
\end{algorithm}
%=====================================================================

%=====================================================================
\section*{Dry Run Example}
Consider the one‑dimensional quadratic
\[
f(w)=(w-3)^2,\qquad \nabla f(w)=2(w-3).
\]
Take $w_0=0$, $\eta_{\max}=0.1$, $M=4$ (so the cosine schedule will be visible after
four steps).  For illustration we ignore stochasticity and set $g_t=\nabla f(w_t)$.

\begin{center}
\begin{tabular}{c|c|c|c|c}
$t$ & $\theta(t)$ & $\eta_t$ (from (1)) & $g_t=2(w_t-3)$ & $w_{t+1}=w_t-\eta_t g_t$\\\hline
0 & $0/4=0$ & $0.1\cdot\frac{1+\cos0}{2}=0.1$ & $2(0-3)=-6$ & $0-0.1(-6)=0.6$\\
1 & $1/4=0.25$ & $0.1\frac{1+\cos(\pi/4)}{2}\approx0.0854$ &
$2(0.6-3)=-4.8$ & $0.6-0.0854(-4.8)\approx1.01$\\
2 & $2/4=0.5$ & $0.1\frac{1+\cos(\pi/2)}{2}=0.05$ &
$2(1.01-3)=-3.98$ & $1.01-0.05(-3.98)\approx1.21$
\end{tabular}
\end{center}
The objective values are $f(0)=9$, $f(0.6)=5.76$, $f(1.01)\approx3.96$,
showing descent while the step size is gradually reduced.

%=====================================================================
\section*{Assumptions}
\begin{assumption}[Smoothness]\label{asm:smooth}
$f$ is $L$‑smooth: for all $x,y\in\mathbb{R}^d$,
\[
\|\nabla f(x)-\nabla f(y)\|\le L\|x-y\|.
\]
Equivalently,
\[
f(y)\le f(x)+\langle\nabla f(x),y-x\rangle+\frac{L}{2}\|y-x\|^2 .
\]
\end{assumption}

\begin{assumption}[Unbiased Stochastic Gradient]\label{asm:unbiased}
For any $t$, conditioned on $w_t$,
\[
\mathbb{E}[g_t\mid w_t]=\nabla f(w_t).
\]
\end{assumption}

\begin{assumption}[Bounded Variance]\label{asm:variance}
There exists $\sigma^2\ge0$ such that
\[
\mathbb{E}\bigl[\|g_t-\nabla f(w_t)\|^2\mid w_t\bigr]\le\sigma^2,
\qquad\forall t.
\]
\end{assumption}

\begin{assumption}[Learning‑rate Schedule]\label{asm:schedule}
The cosine schedule \eqref{1} satisfies
\[
0\le\eta_t\le\eta_{\max},\qquad
\sum_{t=0}^{T-1}\eta_t\ge \frac{2\eta_{\max}}{\pi}T,\qquad
\sum_{t=0}^{T-1}\eta_t^{\,2}\le \frac{\eta_{\max}^{\,2}}{2}T .
\]
\end{assumption}

%=====================================================================
\section*{Main Convergence Theorem}
\begin{theorem}\label{thm:sgdr}
Assume \ref{asm:smooth}–\ref{asm:schedule}.  Let $\{w_t\}$ be generated by
\eqref{2}–\eqref{1}.  Then for any $T\ge1$,
\[
\frac{1}{\sum_{t=0}^{T-1}\eta_t}
\sum_{t=0}^{T-1}\eta_t\,
\mathbb{E}\bigl[\|\nabla f(w_t)\|^2\bigr]
\;\le\;
\frac{2L\bigl(f(w_0)-f^\star\bigr)}{\sum_{t=0}^{T-1}\eta_t}
\;+\;
\frac{L\sigma^2\sum_{t=0}^{T-1}\eta_t^{\,2}}{\sum_{t=0}^{T-1}\eta_t},
\tag{3}
\]
where $f^\star:=\inf_{w}f(w)$.  Consequently, using the bounds in
Assumption~\ref{asm:schedule},
\[
\frac{1}{T}\sum_{t=0}^{T-1}\mathbb{E}\bigl[\|\nabla f(w_t)\|^2\bigr]
\;=\;O\!\left(\frac{1}{\sqrt{T}}\right).
\]
\end{theorem}
%=====================================================================

%=====================================================================
\section*{Theoretical Properties}
\begin{itemize}
  \item \textbf{Stability.}  Because $\eta_t$ is never larger than $\eta_{\max}$,
        the standard SGD stability condition $\eta_{\max}\le 1/L$ guarantees that
        the iterates remain in a bounded region.
  \item \textbf{Hyper‑parameter sensitivity.}
        The period $M$ controls how fast the learning rate decays before a
        restart; larger $M$ yields slower decay and thus a larger
        $\sum_{t}\eta_t$, improving the first term in \eqref{3} but worsening the
        variance term $\sum\eta_t^{2}$.
  \item \textbf{Comparison to plain SGD.}
        Plain SGD with a monotone step‑size $\eta_t=\eta_0/(1+\gamma t)$ also
        achieves $O(1/\sqrt{T})$ on non‑convex smooth objectives.
        SGDR inherits the same rate while often empirically reaching better
        minima because the periodic restarts re‑inject energy into the
        dynamics.
\end{itemize}
%=====================================================================

%=====================================================================
\section*{Preliminaries (Definitions and Lemmas)}
\begin{lemma}[One‑step progress]\label{lem:one-step}
For any $t$,
\[
\mathbb{E}\!\left[f(w_{t+1})\mid w_t\right]
\le
f(w_t)-\eta_t\|\nabla f(w_t)\|^2
+\frac{L\eta_t^{\,2}}{2}\bigl(\|\nabla f(w_t)\|^2+\sigma^2\bigr).
\]
\end{lemma}
\begin{proof}
Condition on $w_t$ and use $L$‑smoothness:
\[
\begin{aligned}
f(w_{t+1})
&\le f(w_t)+\langle\nabla f(w_t),w_{t+1}-w_t\rangle
      +\frac{L}{2}\|w_{t+1}-w_t\|^2 \\[2mm]
&= f(w_t)-\eta_t\langle\nabla f(w_t),g_t\rangle
      +\frac{L\eta_t^{\,2}}{2}\|g_t\|^2 .
\end{aligned}
\]
Take expectation w.r.t. $g_t$ conditioned on $w_t$.
Using Assumptions~\ref{asm:unbiased} and~\ref{asm:variance}:
\[
\mathbb{E}[\langle\nabla f(w_t),g_t\rangle\mid w_t]
 =\|\nabla f(w_t)\|^2,
\qquad
\mathbb{E}[\|g_t\|^2\mid w_t]
 =\|\nabla f(w_t)\|^2+\sigma^2 .
\]
Insert these equalities to obtain the claimed inequality.
\end{proof}
%=====================================================================

%=====================================================================
\section*{Main Proof (Step‑by‑Step Derivation)}
We now prove Theorem~\ref{thm:sgdr}.  All steps are annotated.

\noindent\textbf{[Step 1]}  Take total expectation of Lemma~\ref{lem:one-step}
and sum over $t=0,\dots,T-1$:
\[
\sum_{t=0}^{T-1}\mathbb{E}\!\left[f(w_{t+1})\right]
\le
\sum_{t=0}^{T-1}\mathbb{E}\!\left[f(w_t)\right]
-\sum_{t=0}^{T-1}\eta_t\mathbb{E}\!\left[\|\nabla f(w_t)\|^2\right]
+\frac{L}{2}\sum_{t=0}^{T-1}\eta_t^{\,2}
\mathbb{E}\!\left[\|\nabla f(w_t)\|^2+\sigma^2\right].
\tag{4}
\]

\noindent\textbf{[Step 2]}  The left‑hand side telescopes:
\[
\sum_{t=0}^{T-1}\mathbb{E}[f(w_{t+1})]
=
\mathbb{E}[f(w_T)]\le f(w_0)+\sum_{t=0}^{T-1}\bigl(\mathbb{E}[f(w_{t+1})]-\mathbb{E}[f(w_t)]\bigr).
\]
Subtract $\sum_{t=0}^{T-1}\mathbb{E}[f(w_t)]$ from both sides of \eqref{4} to obtain
\[
\mathbb{E}[f(w_T)]-f(w_0)
\le
-\sum_{t=0}^{T-1}\eta_t\mathbb{E}\!\left[\|\nabla f(w_t)\|^2\right]
+\frac{L}{2}\sum_{t=0}^{T-1}\eta_t^{\,2}
\mathbb{E}\!\left[\|\nabla f(w_t)\|^2\right]
+\frac{L\sigma^2}{2}\sum_{t=0}^{T-1}\eta_t^{\,2}.
\tag{5}
\]

\noindent\textbf{[Step 3]}  Rearrange \eqref{5} and bring the terms involving
$\mathbb{E}\|\nabla f(w_t)\|^2$ to the left:
\[
\sum_{t=0}^{T-1}\Bigl(\eta_t-\frac{L}{2}\eta_t^{\,2}\Bigr)
\mathbb{E}\!\left[\|\nabla f(w_t)\|^2\right]
\le
f(w_0)-\mathbb{E}[f(w_T)]
+\frac{L\sigma^2}{2}\sum_{t=0}^{T-1}\eta_t^{\,2}.
\tag{6}
\]

\noindent\textbf{[Step 4]}  Since $f(w_T)\ge f^\star$, we can replace
$f(w_0)-\mathbb{E}[f(w_T)]\le f(w_0)-f^\star$.
Moreover, for any $t$ we have $\eta_t\le\eta_{\max}\le 1/L$ (standard stability
requirement).  Hence $\eta_t-\frac{L}{2}\eta_t^{\,2}
\ge\frac{\eta_t}{2}$.
Thus
\[
\frac{1}{2}\sum_{t=0}^{T-1}\eta_t
\mathbb{E}\!\left[\|\nabla f(w_t)\|^2\right]
\le
f(w_0)-f^\star
+\frac{L\sigma^2}{2}\sum_{t=0}^{T-1}\eta_t^{\,2}.
\tag{7}
\]

\noindent\textbf{[Step 5]}  Divide both sides of \eqref{7} by
$\frac12\sum_{t=0}^{T-1}\eta_t$ and obtain exactly the bound \eqref{3}:
\[
\frac{\sum_{t=0}^{T-1}\eta_t\,
\mathbb{E}[\|\nabla f(w_t)\|^2]}
{\sum_{t=0}^{T-1}\eta_t}
\le
\frac{2L\bigl(f(w_0)-f^\star\bigr)}{\sum_{t=0}^{T-1}\eta_t}
+\frac{L\sigma^2\sum_{t=0}^{T-1}\eta_t^{\,2}}
      {\sum_{t=0}^{T-1}\eta_t}.
\]
\qed

%=====================================================================
\section*{Rate Derivation (With Constants)}
Using the explicit bounds from Assumption~\ref{asm:schedule},
\[
\sum_{t=0}^{T-1}\eta_t\;\ge\;\frac{2\eta_{\max}}{\pi}\,T,
\qquad
\sum_{t=0}^{T-1}\eta_t^{\,2}\;\le\;\frac{\eta_{\max}^{\,2}}{2}\,T .
\]
Insert them into \eqref{3}:
\[
\frac{1}{T}\sum_{t=0}^{T-1}\mathbb{E}\!\left[\|\nabla f(w_t)\|^2\right]
\le
\underbrace{\frac{\pi L\bigl(f(w_0)-f^\star\bigr)}{\eta_{\max}}}_{\displaystyle
C_1}\frac{1}{T}
\;+\;
\underbrace{\frac{\pi L\sigma^2\eta_{\max}}{4}}_{\displaystyle C_2}
\frac{1}{\sqrt{T}} .
\]
Since the $1/\sqrt{T}$ term dominates for large $T$, we obtain the
asymptotic rate $O(1/\sqrt{T})$.

%=====================================================================
\section*{Special Cases}
\begin{itemize}
  \item \textbf{Deterministic gradients ($\sigma=0$).}
        The variance term disappears and the bound simplifies to
        $O(1/T)$, matching the classic result for gradient descent with a
        diminishing step size.
  \item \textbf{Constant step size $\eta_t=\eta_{\max}$.}
        The schedule reduces to plain SGD; the bound becomes
        $O(1/T)$ for the first term but the variance term scales as
        $O(1)$, yielding the well‑known $O(1/\sqrt{T})$ stationary‑point
        rate.
  \item \textbf{Infinite horizon ($T\to\infty$).}
        The average gradient norm converges to $0$ because the right‑hand
        side of \eqref{3} tends to $0$ as $\sum\eta_t\to\infty$ while
        $\sum\eta_t^{\,2}=o(\sum\eta_t)$ for the cosine schedule.
\end{itemize}
%=====================================================================

\end{document}